\chapter{ML-модели для операционного и финансового управления}

\section{Зачем нужны прогнозные модели}

Финансовый аудит -- это взгляд назад. Он показывает, что произошло, где возникли кассовые провалы, какие поставщики занимают критическую долю расходов. Но для управленца не менее важен взгляд вперёд: что произойдёт на следующей неделе, когда лучше закупить сырьё, в каком месяце стоит подготовить резерв ликвидности.

В проекте разработаны несколько прогнозных моделей, каждая из которых решает отдельную управленческую задачу: прогноз спроса на товары, прогноз объёма возвратов и прогноз дневного cash flow с детектором аномалий. Дополнительно реализована система мониторинга качества моделей -- чтобы было понятно, когда модель начинает <<врать>> и пора её переобучить.

\section{Прогноз спроса}

\subsection{Постановка задачи}

Задача прогноза спроса -- предсказать, сколько единиц каждого товара будет продано на следующей неделе. Для производственной компании это напрямую влияет на объём производства: слишком много -- товар остаётся и портится (кондитерские изделия имеют ограниченный срок годности), слишком мало -- недополученные продажи и недовольство партнёров.

\subsection{Данные и признаки}

Обучающая выборка -- продажи по контрагентам за 14 месяцев (только реализации, без корректировок). Гранулярность: неделя × товар. При такой истории классические модели временных рядов ARIMA на каждый товар отдельно дадут нестабильные результаты -- слишком мало точек. Gradient boosting позволяет объединить данные по всем товарам в одну модель и за счёт категориальных признаков разделять закономерности между позициями и сетями.

Используемые признаки:
\begin{itemize}
  \item лаги объёма продаж: прошлая неделя, две недели назад, четыре недели назад;
  \item скользящие суммы за 4 и 8 недель (тренд и сезонность);
  \item календарные признаки: номер недели, месяц, флаг предпраздничных недель;
  \item категориальные: наименование товара, торговая сеть, юридическое лицо покупателя.
\end{itemize}

Для сезонности кондитерских изделий предпраздничные недели (Новый год, 8 марта, майские) особенно важны: спрос в эти периоды резко отличается от средней недели.

\subsection{Бейзлайн и модель}

В качестве бейзлайнов использованы: наивный прогноз (<<продажи будут как на прошлой неделе>>) и скользящая средняя за 4 недели. Градиентный бустинг (CatBoost) сравнивается с бейзлайнами по метрике WAPE:

\[
\text{WAPE} = \frac{\sum |y_i - \hat{y}_i|}{\sum y_i}
\]

где $y_i$ -- фактические продажи, $\hat{y}_i$ -- прогноз. WAPE предпочтительнее MAPE для данной задачи, поскольку не штрафует за ошибки на малых объёмах непропорционально.

\section{Прогноз возвратов}

\subsection{Природа возвратов}

Возвраты кондитерских изделий связаны в первую очередь со сроком годности: товар, не успевший продаться, возвращается производителю. Объём возвратов за анализируемый период составил около 8,3 млн руб. от 53,3 млн руб. валовых продаж, то есть около 15,6\% -- это существенная цифра, которую нельзя игнорировать при планировании производства.

\subsection{Модель и применение}

Модель прогноза возвратов -- регрессия объёма корректировок реализации на следующей неделе по товару и сети. Ключевой признак -- исторический return rate этого товара в этой сети: если в прошлые периоды возвращалось 20\% отгруженного, это сильный сигнал для следующего прогноза. Дополнительно вводится \texttt{срок\_годности\_дней} из справочника товаров: изделия с коротким сроком имеют системно более высокий риск возврата.

\begin{center}
\begin{tabular}{lrr}
\toprule
\textbf{Модель} & \textbf{CV WAPE} & \textbf{Recent WAPE} \\
\midrule
Возвраты (catboost\_mae) & 0,709 & 0,668 \\
\bottomrule
\end{tabular}
\end{center}

Качество модели возвратов ниже, чем у cash flow, -- это ожидаемо. Корректировки реализации по природе своей шумные: один крупный возврат от одного клиента может в несколько раз превысить среднее. Тем не менее модель полезна как риск-фактор: она показывает ожидаемый объём возврата, который вычитается из рекомендации производства.

\section{Рекомендация объёма производства}

Рекомендация объёма производства -- это не отдельная ML-модель, а правило поверх прогнозов спроса и возвратов:

\[
\text{Рекомендация} = \frac{\text{Прогноз спроса} - \text{Остаток на складе}}{1 - \text{Ожидаемый return rate}}
\]

Логика проста: нужно произвести достаточно, чтобы покрыть ожидаемые продажи за вычетом того, что уже лежит на складе, с поправкой на то, что часть отгруженного вернётся. При ограниченных мощностях дополнительно учитывается маржинальность товара: производить сначала то, что даёт больший вклад на единицу.

Такой компонент намеренно оставлен детерминированным -- без отдельного ML-обучения. При 16 товарных позициях и 14 месяцах истории сложная модель здесь переусложнение; прозрачное правило лучше поддаётся проверке и коррекции.

\section{Прогноз дневного cash flow}

\subsection{Особенность задачи}

Ежедневный денежный поток производственной компании -- это так называемый \textit{intermittent time series}: многие дни нулевые (нет ни поступлений, ни крупных списаний), а отдельные дни дают большие всплески. Классическая регрессия плохо работает с такими рядами -- она предсказывает <<среднее>> там, где нужно уловить <<нулевые дни vs дни с платежом>>.

\subsection{Hurdle-подход}

Для решения этой проблемы применяется hurdle-модель: двухэтапный подход, где сначала оценивается вероятность ненулевого дня (бинарная классификация), а затем -- сумма при условии, что день ненулевой (регрессия):

\[
\hat{y}_t = P(\text{день ненулевой}_t) \times \hat{y}^+_t
\]

Это позволяет корректно обрабатывать недели без крупных поступлений и дни с неожиданно большими платежами.

\subsection{Результаты}

\begin{center}
\begin{tabular}{lrr}
\toprule
\textbf{Модель} & \textbf{CV WAPE} & \textbf{Recent WAPE} \\
\midrule
Cash Flow: поступления (hurdle\_catboost) & 0,237 & 0,122 \\
Cash Flow: списания (hurdle\_catboost) & 0,340 & 0,220 \\
\bottomrule
\end{tabular}
\end{center}

Recent WAPE 0,122 для поступлений -- это рабочее качество для управленческого предупреждения: на горизонте 30 дней модель даёт разумную оценку ожидаемых денежных дней и примерных сумм. Для списаний ошибка выше (0,220) -- это связано с тем, что крупные разовые платежи поставщикам создают непредсказуемые всплески. Для улучшения этой части модели в продовой версии желательно добавить явный календарь обязательных платежей (аренда, зарплата, налоги, кредиты) как внешние регрессоры.

\section{Детектор аномалий поступлений}

Помимо прогнозирования агрегированного потока, в проекте реализован детектор аномалий. Его задача -- другая: найти конкретную операцию, которая выбивается из типичного паттерна данного контрагента.

Модель работает по принципу unsupervised learning: статистические пороги на основе z-score (насколько текущая сумма отличается от среднего по истории этого контрагента) и EWMA. Для финансового контроля это предпочтительнее сложных нейросетевых методов: результат прозрачен, легко объясним и не требует размеченных примеров <<аномалии>>.

В текущем прогоне детектор выявил \textbf{42 транзакции} на дополнительную проверку и \textbf{3 аномальных дня} за последние 30 дней. Результат детектора не блокирует платежи автоматически -- он формирует очередь на ручной контроль с указанием причины (нетипичная сумма, новый контрагент, нестандартный день платежа).

\section{Мониторинг качества и разладка}

\subsection{Проблема модельного дрейфа}

Любая модель, обученная на исторических данных, постепенно устаревает: меняются цены, сезонность, клиенты, условия работы. Если не отслеживать качество модели на свежих данных, можно долго использовать прогнозы, которые уже не отражают реальность.

\subsection{Система мониторинга}

В проекте реализована система сравнения свежей ошибки с ошибкой на кросс-валидации и расчёта PSI (Population Stability Index) по распределениям cash flow:

\[
\text{PSI} = \sum_{i} \left(p_i - q_i\right) \ln\frac{p_i}{q_i}
\]

где $p_i$ -- доля в исходном обучающем распределении, $q_i$ -- доля в новых данных. PSI меньше 0,10 -- распределение стабильно; от 0,10 до 0,25 -- требует внимания; выше 0,25 -- распределение значимо изменилось.

По итогам статусы делятся на три группы:
\begin{itemize}
  \item \textbf{ok} -- свежая ошибка не сильно хуже CV-ошибки, PSI в норме;
  \item \textbf{watch} -- небольшое ухудшение; модель публикуется, но требует наблюдения;
  \item \textbf{retrain} -- значимое ухудшение; модель нужно переобучить и отправить на ревью перед публикацией.
\end{itemize}

\subsection{Текущий статус}

Текущий статус всех моделей -- \textbf{retrain}. Свежая ошибка моделей при этом остаётся лучше CV-ошибки (то есть модели не деградировали по качеству), но распределения cash flow существенно изменились: PSI по поступлениям -- 1,032, PSI по списаниям -- 1,165. Это намного выше порога 0,25 -- распределения радикально сдвинулись.

Это типичная ситуация при первом запуске в новых данных: исторические данные для обучения и свежие операции могут различаться просто из-за сезонных паттернов или разовых событий. Система корректно фиксирует факт разладки и рекомендует переобучение -- это именно то поведение, которое нужно для production-loop.

\subsection{Production-loop переобучения}

Пайплайн мониторинга устроен как контур регулярного обновления:
\begin{enumerate}
  \item загрузить свежие данные из итоговых таблиц;
  \item пересчитать признаки без утечки будущего;
  \item переобучить модели на всей доступной истории;
  \item сравнить свежую ошибку с CV-ошибкой и посчитать PSI;
  \item при статусе \texttt{ok} или \texttt{watch} -- публиковать прогноз в дашборд;
  \item при статусе \texttt{retrain} -- отправить результат на ревью перед публикацией.
\end{enumerate}

Такой подход решает одну из ключевых проблем ML в малом бизнесе: модели не просто обучаются один раз, они встроены в регулярный управленческий контур и адаптируются по мере поступления новых данных.

\section{Ограничения моделей}

Чтобы результаты работы не воспринимались как готовое production-решение, важно честно зафиксировать ограничения.

\textbf{Короткая история.} 14 месяцев данных -- это мало для моделей временных рядов. Годовые сезонные паттерны просматриваются только один раз, что делает их оценки менее надёжными.

\textbf{Чувствительность к разовым событиям.} Крупный разовый закупочный платёж (например, авансовая оплата большой партии) сильно искажает распределение cash flow. Без явного календаря обязательных платежей модель списаний плохо предсказывает такие всплески.

\textbf{Отсутствие фактических остатков на счёте.} Модель прогнозирует изменение остатка (дельту), а не абсолютный баланс. Для практического применения нужна ежедневная выгрузка фактического остатка.

\textbf{Неполный справочник товаров.} Рекомендация объёма производства требует данных о сроке годности и себестоимости по каждой позиции. Без заполненного справочника этот компонент работает частично.

Несмотря на эти ограничения, модели уже сейчас пригодны для управленческого использования в режиме <<раннего предупреждения>>: они показывают, в каком направлении движется денежный поток, а не претендуют на точность до рубля.
